\section{آماده‌کردن الگو}

الگوی برش شکل و اندازۀ قطعۀ خام را مشخص می‌کند. الگو ممکن است خط دوخت، خط تا و علامت‌های تطبیق هم داشته باشد. برای خواندن آن، ابتدا نام قطعه، تعداد، ضخامت و سمت رو یا پشت را مشخص کنید. در الگوی چندقطعه‌ای، قطعات هم‌جفت و آینه‌ای را پیش از مصرف چرم روی کاغذ کنار هم بگذارید و طول لبه‌های مشترک را بسنجید.

\subsection{خط برش، خط دوخت و خط تا}

در شکل~\ref{fig:lines}، خط پیوستۀ \lr{A} مرز برش، خط‌چین \lr{B} راهنمای دوخت، خط نقطه‌چین \lr{C} راهنمای تا و دایرۀ \lr{D} محل سوراخ است. این قرارداد برای همین جزوه است. راهنمای هر الگو تعیین می‌کند که دایره فقط مرکز را نشان می‌دهد یا مرز سوراخ را هم مشخص می‌کند.

\begin{figure}[!htbp]
  \centering
  \PatternLinesDiagram
  \caption{نمونۀ ترسیمی چهار علامت روی الگو.}
  \label{fig:lines}
\end{figure}

جای دوخت می‌تواند به فاصلۀ خط دوخت تا لبۀ برش اشاره کند. در الگوی آماده، بررسی کنید که این فاصله در شکل لحاظ شده باشد. فاصله را در گوشه‌ها یا دور منحنی هم بسنجید. خط تا روی الگو محل اسمی تا را نشان می‌دهد. وقتی چند لایه روی هم قرار می‌گیرند، برای تعیین طول مفید و محل تا، تکه‌ای با همان لایه‌بندی بسازید.

اگر با \lr{SVG} یا \lr{DXF} کار می‌کنید، واحد را میلی‌متر بگیرید و دست‌کم یک طول معلوم را پس از بازکردن فایل در نرم‌افزار مقصد اندازه بگیرید. برای ساخت، مسیرهای برش، علامت‌های دوخت و خطوط تا را روی لایه‌های جدا بگذارید. محیط بیرونی و سوراخ‌های برش باید مسیرهای بسته داشته باشند. درز باز، خط تکراری و تقاطع ناخواسته را پیش از سفارش شابلون یا قالب رفع کنید.

\subsection{چاپ تک‌برگ و چندبرگ}

صفحۀ تک‌برگ یا برگه‌های از پیش صفحه‌بندی‌شده را با \lr{Actual Size} یا مقیاس ۱۰۰ درصد روی اندازۀ کاغذ تعیین‌شده چاپ کنید. حالت \lr{Fit} مقیاس را تغییر می‌دهد. خط کنترل را پس از چاپ از مرکز نشانۀ ابتدا تا مرکز نشانۀ انتها اندازه بگیرید. در الگوی دقیق، هر دو جهت کاغذ را کنترل کنید. اگر لبه‌ای از شکل بیرون حاشیۀ قابل‌چاپ پرینتر افتاده است، طرح صفحه یا دستگاه را بررسی کنید \cite{adobe-size}.

اگر فایل در یک صفحۀ بزرگ طراحی شده، چاپ \lr{Poster} امکان تقسیم آن به چند برگ را می‌دهد. در این حالت، مقیاس تقسیم، هم‌پوشانی و علامت‌های برش را در پیش‌نمایش و سپس روی یک برگ آزمایشی بررسی کنید. برای الگوی چندبرگی همین جزوه، هر صفحه از قبل یک برگ \lr{A4} است و به تقسیم دوباره نیاز ندارد \cite{adobe-poster}.

چندبرگی‌کردن فقط جای الگو روی کاغذ را تغییر می‌دهد. تغییر ابعاد محصول، تصمیم طراحی است. مثلاً با بزرگ‌کردن تمام صفحه به ۱۱۰ درصد، فاصلۀ دوخت از لبه و قطر سوراخ هم بزرگ می‌شوند. اگر کارت، زیپ یا یراق همان اندازه قبلی را دارند، ابتدا ابعاد قطعه را تغییر دهید و جای دوخت و سوراخ‌ها را جداگانه بازبینی کنید.

\section{شابلون، چیدمان و علامت‌گذاری}

شابلون الگوی فیزیکیِ قرارگرفته روی چرم است. پیش از ساخت شابلون بادوام، نسخه‌ای کاغذی را برش دهید و قطعات را موقتاً مونتاژ کنید. پس از تأیید تناسب، خط مرجع، نام قطعه و تعداد را روی شابلون ثبت کنید. اندازۀ کاغذ یا شابلون را دوباره با طول اصلی مقایسه کنید.

\begin{table}[!htbp]
  \centering
  \caption{انتخاب شابلون بر اساس شکل قطعه و تعداد برش}
  \label{tab:template}
  \begin{tabular}{@{}p{29mm}p{50mm}p{67mm}@{}}
    \toprule
    جنس & کاربرد مناسب & نکته در انتقال \\
    \midrule
    کاغذ & نمونه‌سازی، الگوی بزرگ یا طرحی که هنوز اصلاح می‌شود & کاغذ را با وزنه ثابت کنید. خط را با فشار کم منتقل کنید تا کاغذ جابه‌جا نشود. \\
    مقوای نازک & برش تکراری در تعداد کم & در گوشه‌ها تغییر شکل لبه یا لب‌پرشدن را بررسی کنید. \\
    اکریلیک & برش‌های تکراری و شابلون شفافی که جای عیب سطح را نشان دهد & ضخامت، لغزش و دسترسی نوک علامت‌زن به گوشه‌ها را بررسی کنید. برای هدایت تیغ از راهنمای مناسب ابزار استفاده کنید. \\
    \bottomrule
  \end{tabular}
\end{table}

جدول~\ref{tab:template} گزینه‌های شروع کار را مقایسه می‌کند. برای نقطه‌های داخلی، انتقال مرکز با علامت‌زن معمولاً دقیق‌تر از دنبال‌کردن یک سوراخ بزرگ در شابلون است. رو و پشت قطعه، تعداد و جهت کشش را روی شابلون مشخص کنید. کد قطعاتی را که شبیه هم‌اند، پیش از برداشتن شابلون به پشت چرم منتقل کنید.

\subsection{چیدن قطعات روی پوست}

چرم را صاف و در حالت طبیعی روی میز بگذارید. کشیدن چرم نرم هنگام علامت‌گذاری ممکن است شکل قطعه را پس از آزادشدن تغییر دهد. ابتدا قطعات بزرگ و بندهای بلند را در ناحیۀ قابل‌استفاده جا بدهید و قطعات کوچک را در فواصل باقی‌مانده بچینید. پاره، خراش عمیق، ضخامت نامنظم و کشش نامناسب را به‌ویژه در محل دوخت، خط تا و اتصال یراق در نظر بگیرید.

برای دو قطعۀ آینه‌ای، شابلون را روی همان سمت چرم یک بار برگردانید یا دو شابلون مشخصِ چپ و راست بسازید. پیش از برش، هر دو قطعه را از سمت رویه در کنار هم تصور کنید. در سطح پرزدار، همۀ قطعاتی که کنار هم دیده می‌شوند باید جهت پرز موردنظر را داشته باشند. برای قطعۀ باربر، جهت مناسب را با آزمون روی تکه‌ای از همان ناحیه تعیین کنید.

\subsection{انتقال خطوط و علامت‌ها}

ابزار علامت‌گذاری را روی دورریز همان چرم امتحان کنید. برای رویۀ حساس، می‌توان خط را روی پشت یا بیرون از مرز برش منتقل کرد. اگر از پشت علامت می‌زنید، جهت قطعات نامتقارن را دوباره کنترل کنید. امتداد خط‌کش را با لبه‌های شابلون تطبیق دهید و علامت‌ها را بدون کشیدن چرم بزنید. اگر خط‌زن اثری با پهنا ایجاد می‌کند، پیش از برش تعیین کنید مرکزِ خط یا یکی از لبه‌های آن مرجع است.

خط دوخت، مرکز پانچ و محدودۀ لویس را با علامت‌های متمایز منتقل کنید. برش روی خطِ اشتباه را قبل از شروع با شمارش قطعات، مقایسۀ علامت‌ها با شکل~\ref{fig:lines} و کنترل چند فاصلۀ کلیدی پیشگیری کنید.
